\section{HDI and ETI: Credible Interval Variants}\label{app:hdi}

Given a posterior density $p(\theta \mid \mathbf{x})$, a $(1-\alpha)$ credible interval
$[L, U]$ satisfies
\begin{equation}\label{eq:credible_interval}
  \int_L^U p(\theta \mid \mathbf{x})\,d\theta = 1 - \alpha.
\end{equation}
Infinitely many intervals satisfy this constraint. This paper uses the Highest Density
Interval throughout; the Equal-Tailed Interval is described below for comparison.

\subsection*{Highest Density Interval (HDI)}

The HDI is the shortest interval satisfying Equation~\ref{eq:credible_interval}:
\begin{equation}\label{eq:hdi}
  [L^*, U^*] = \underset{[L,U]:\;\int_L^U p(\theta|\mathbf{x})\,d\theta\,=\,1-\alpha}{\arg\min}\;(U - L).
\end{equation}
An equivalent characterisation for unimodal posteriors: $[L^*, U^*]$ is the unique
interval for which $p(L^* \mid \mathbf{x}) = p(U^* \mid \mathbf{x})$, so every point
inside has posterior density at least as high as every point outside (hence ``highest
density'').\footnote{For multimodal posteriors, multiple intervals can satisfy the
equal-density boundary condition, and the HDI may be a union of disjoint intervals;
the shortest-interval definition in Equation~\ref{eq:hdi} selects the one with minimum
total width. All posteriors in this paper are unimodal.}

\subsection*{Equal-Tailed Interval (ETI)}

The ETI places $\alpha/2$ probability mass in each tail:
\begin{equation}\label{eq:eti}
  L_{\rm ETI} = F^{-1}\!\left(\tfrac{\alpha}{2}\right), \qquad
  U_{\rm ETI} = F^{-1}\!\left(1 - \tfrac{\alpha}{2}\right),
\end{equation}
where $F^{-1}$ is the posterior quantile function. The ETI is simple to compute and
uniquely defined, but it is not generally the shortest interval.
For symmetric, unimodal posteriors the HDI and ETI coincide; they diverge for skewed
distributions --- for instance, a Beta posterior near $\theta = 0$ or $\theta = 1$ ---
where the HDI is noticeably shorter and therefore preferred.

\subsection*{Numerical Computation}

The HDI is computed numerically by minimising the interval width over the lower tail
probability $\ell$ using unconstrained Nelder-Mead optimisation (\texttt{HDIofICDF}
in the accompanying code, implemented via \texttt{scipy.optimize.fmin}):
\begin{equation}
  \min_{\ell \,\geq\, 0}\left[F^{-1}(1-\alpha+\ell) - F^{-1}(\ell)\right].
\end{equation}
The minimising $\ell^*$ gives $L^* = F^{-1}(\ell^*)$ and
$U^* = F^{-1}(1-\alpha+\ell^*)$.
For the Binomial model used in this work, under a flat (uniform) prior
$\mathrm{Beta}(1,1)$, the posterior after $s$ successes and $f$ failures is
$\mathrm{Beta}(s+1,\,f+1)$, and the HDI is found by applying the above minimisation
to the Beta quantile function.
% Prior convention note (for the author):
% The flat prior Beta(1,1) adds one pseudo-count to both successes and failures,
% giving posterior Beta(s+1, f+1). This is the standard Bayesian conjugate update
% for Bernoulli data and is the convention used by Kruschke (DBDA2).
%
% The alternative is the Haldane prior Beta(0,0), named after geneticist J.B.S.
% Haldane. It is an improper prior (it does not integrate to 1) that assigns infinite
% weight to the boundary values theta=0 and theta=1. Under this convention the
% posterior is simply Beta(s, f) --- i.e. the raw count data directly parameterise
% the posterior, with no pseudo-counts. This can feel natural in code (just pass
% successes and failures straight to scipy.stats.beta), but it is undefined when
% s=0 or f=0, requiring an ad hoc +1 workaround for degenerate cases. The earlier
% version of this code used Beta(s, f), implicitly adopting the Haldane convention
% without stating it. For the sample sizes used in this paper (N >> 1) the numerical
% difference between Beta(s,f) and Beta(s+1,f+1) is negligible, but the flat prior
% is the better-motivated and more transparent default.
%
% CODE NOTE: switching to Beta(s+1, f+1) here requires updating
% successes_failures_to_hdi_ci_limits() calls in utils_stats.py to pass
% (s+1, f+1) instead of (s, f). The edge-case +1 workarounds in
% stop_decision_multiple_experiments_pitg() can then be removed.

The same \texttt{HDIofICDF} function handles the continuous setting by substituting
the Student-$t$ quantile function with location equal to the sample mean and scale
equal to the standard error (\texttt{continuous\_hdi\_ci\_limits} in the accompanying
code); the corresponding sample-size formula is derived in
Appendix~\ref{app:expected_stop}.
